\section{Introduction}\label{sec:intro}

\subsection{Bounding cochains in the pillowcase}
Let $K\subset S^3$ be a knot. On the gauge-theoretic side sits the reduced singular instanton knot
homology $\Inat(K)$ of Kronheimer and Mrowka~\cite{KM1,KM2}. On the symplectic side sits the pillowcase
program of Hedden, Herald and Kirk~\cite{HHK1,HHK2}. Cut $K$ along a Conway sphere into two $2$-tangles.
The traceless $\SU$ character variety of the four-punctured sphere is the \emph{pillowcase}, a two-sphere
with four orbifold points, and after a holonomy perturbation each tangle determines an immersed curve in
it; the tangle that carries the earring of the reduced theory gives a figure-eight curve. Hedden, Herald
and Kirk conjecture that for a suitable decomposition the Lagrangian Floer homology of the two curves is
isomorphic to $\Inat(K;\F_2)$~\cite[Conj.~6.5]{HHK2}. For general decompositions figure-eight bubbles at
the self-intersections of the curves deform the differential. Cazassus, Herald, Kirk and Kotelskiy
conjecture that each tangle is assigned an immersed curve together with a bounding cochain, and that the
deformed theory recovers $\Inat(K)$~\cite[Conj.~D]{CHKK}. Conditional on this conjecture, Smith proves that
the cochain on one Conway-sum tangle in $P(-2,3,5)$ cannot vanish~\cite[proof of Thm.~5.9]{Smith}.

This paper studies Smith's decomposition of the pretzel knots $K_q=P(-2,3,q)$, $q\ge3$
odd~\eqref{eq:family}, as the numerator closure of a sum of three rational tangles,
$K_q=\operatorname{num}(Q_{-1/2}+Q_{1/3}+Q_{1/q})$, with the earring on the $-2$ band. Smith shows that
the Floer homology with zero bounding cochain has the wrong rank at
$q=5$~\cite[proof of Thm.~5.9]{Smith}, and at $q=7$ the same holds for a piecewise-linear model of his curve
(Proposition~\ref{prop:q7-typeD}). We work over the algebra $\mathcal B$ of Kotelskiy, Watson and
Zibrowius~\cite{KWZ}, the endomorphism algebra of two arcs in the four-punctured sphere. For $q=5,7$ we
construct Maurer--Cartan elements for the twisted complex of the model over $\mathcal B$ that correct the
rank, and we identify the resulting objects with the Bar-Natan invariant of the tangle $Q_{1/3}+Q_{1/q}$ up
to one further term (Theorems~\ref{thm:structure} and~\ref{thm:pairing}). The companion paper~\cite{WuebbenII} computes
$\Inat(K_q;\Z)\cong\Z^{q+2}$ for every odd $q\ge3$~\cite[Thm.~1.1]{WuebbenII}; these instanton dimensions
are the targets of the pairing dimensions computed here. It also shows that for a different decomposition
of $K_q$, obtained from the tangle of Hedden, Herald and Kirk for $T(3,5)$ by regluing, the pillowcase Floer
homology with zero bounding cochain is isomorphic to $\Inat(K_q;\F_2)$ as a $\Z/4$-graded vector space, for
odd $q\ge7$~\cite[Thm.~1.2]{WuebbenII}.

\subsection{Smith's decomposition}
In Smith's decomposition the earring is placed on the summand $Q_{-1/2}$, which gives the tangle
$\Qh{-1/2}$; for a slope $r$ we write $\Qh{r}$ for the tangle $Q_r$ with an earring, whose curve $E_r$
surrounds the line of $Q_{-r}$ (\S\S\ref{ss:models} and~\ref{ss:kh-smith}). The other tangle is
$T_q=Q_{1/3}+Q_{1/q}$. We work in Smith's coordinates on the pillowcase, with his holonomy perturbation
for small negative values of its parameter $t$ (\S\ref{ss:smith-chart}). Following Cazassus, Herald, Kirk
and Kotelskiy~\cite[\S11]{CHKK}, we encode curves as twisted complexes over $\mathcal B$. The higher
products of $\mathcal B$ vanish, so the Maurer--Cartan equation for a deformation $b$ of a twisted complex
$(V,\delta)$ is the finite identity $(\delta+b)^2=0$. The pairing dimension of two twisted complexes is the
dimension of the homology of their morphism complex, a Floer homology dimension for curves in the position
required by~\cite[Thm.~5.22]{KWZ} (\S\ref{ss:kwz}). We write $N_q$ for the twisted complex of a
piecewise-linear model of Smith's curve for $T_q$, and $\mathcal E_r$ for that of the earring curve $E_r$,
with $\mathcal E=\mathcal E_{-1/2}$.

Under the equivalence of $\operatorname{Tw}(\mathcal B)$ with the twisted complexes over the relevant
wrapped Fukaya category of the pillowcase~\cite[\S11.5]{CHKK}, the deformed complexes $(N_q,b)$ play the role
of immersed curves with bounding cochains. The switch of a smoothing at a self-intersection, defined below,
is the twisted-complex form of the Lagrangian surgery there, which in immersed Floer theory corresponds to a degree-one bounding cochain at that point (for curves on surfaces compare Auroux and Smith, \S\ref{ss:disc-prior}). We do not express the switches in the immersed Floer complex
of Akaho and Joyce~\cite{AkahoJoyce} (Remark~\ref{rem:q7-scope}).

For $q=7$ the pairing dimension of $\mathcal E$ with $N_7$ is $7$, not $9=\dim\Inat(K_7)$
(Proposition~\ref{prop:q7-typeD}). Each self-intersection of the model lifts to two double points of its lift to the torus double cover, exchanged by the elliptic involution; we call the pair a \emph{self-intersection orbit}. A smoothing of the
model at a self-intersection orbit that keeps the generators of $N_7$ changes its differential by a matrix
$b_s$, the \emph{switch} of the smoothing. At four self-intersection orbits, labelled
$S_{18},S_{25},S_{69},S_{74}$ in \S\ref{ss:smoothings}, smoothing gives a switch $b_s$ that is a
Maurer--Cartan element and raises the pairing dimension to $9$, and the four objects $(N_7,b_s)$ fall into
three homotopy-equivalence classes (Corollary~\ref{cor:q7-classes}). Replacing $\Qh{-1/2}$ by $\Qh{-3/4}$
closes $T_7$ into a second knot $K_*$. Its reduced Khovanov homology is free of rank $31$, which equals the
coefficient norm of its Alexander polynomial, so the bounds of Kronheimer and Mrowka~\cite{KMAlex,KM2} give
$\Inat(K_*;\Z)\cong\Z^{31}$, and the pairing dimensions with the second earring separate the three classes
(Proposition~\ref{prop:aux-closure}).

\emph{Conditional selection at $q=7$.} Assume the conjecture of Cazassus, Herald, Kirk and Kotelskiy, and
assume Hypothesis~\ref{hyp:q7-support}: the assigned object is homotopy equivalent to the oriented smoothing
of the model at a single self-intersection orbit, or to the deformation of $N_7$ by a combination of the four
switches above.
Then the object assigned to $T_7$ is homotopy equivalent to the common class of $(N_7,b_{S_{18}})$ and
$(N_7,b_{S_{25}})$ (Corollary~\ref{cor:aux-conditional}). The orbits $S_{18}$ and $S_{25}$ are the two at
which a double point cuts the lift into two loops the shorter of which has torus class $(2,2)$ up to sign
(Proposition~\ref{prop:aux-selection}). For positive $t$ two of the four smoothings give $5$ and $7$ instead
of $9$, and the two-closure selection holds for both signs (Remark~\ref{rem:sign-pairings}). We also show that Gao's wrapped representability
theorem~\cite{Gao} cannot be applied to the Lagrangian correspondence induced by the line of $Q_{1/3}$,
which is immersed with a triple point (Proposition~\ref{prop:partial-triple}).

Kotelskiy, Watson and Zibrowius assign to $T_q$ its Bar-Natan invariant $\widetilde{\mathrm{BN}}(T_q)$, an
immersed multicurve with local systems in the four-punctured sphere~\cite{KWZ}. Drawn in Smith's
coordinates it becomes an object $\mathrm{BN}_q$ of $\operatorname{Tw}(\mathcal B)$ (\S\ref{ss:kh-smith}).
For $r\in\Q\cup\{\infty\}$ let $K_r(q)=\operatorname{num}(Q_r+T_q)$, so that $K_{-1/2}(q)=K_q$. A compact
curve in the pillowcase is \emph{linear} if its lifts to the plane lie close to parallel straight lines and
are graphs over them (Definition~\ref{def:linear}). Two twisted complexes are \emph{strictly isomorphic} if
a bijection of their generators that preserves idempotents carries one differential to the other entry by
entry (\S\ref{ss:kwz}). For odd $q$ put $r_q=(1-q)/(3q-4)$. Let $\alpha_q$ be the straight arc between
$(0,0)$ and $(\pi,0)$ that the earring curve $E_{r_q}$ surrounds, and let $c$ be the boundary of a regular
neighborhood of the straight arc that $E_{-1}$ surrounds; $c$ separates $\{(\pi,0),(0,\pi)\}$ from
$\{(0,0),(\pi,\pi)\}$ (\S\ref{ss:kh-structure}). For $q\in\{5,7\}$ let $b$ be one of the selected elements
$b_{S_{18}},b_{S_{25}}$ at $q=7$, or at $q=5$ the switch $b_5$ of a smoothing that raises the pairing
dimension with $\mathcal E$ from $5$ to $7$ (\S\ref{ss:kh-structure}).

\begin{theorem}\label{thm:structure}
Let $q\in\{5,7\}$.
\begin{enumerate}
\item The object $(N_q,b)$ is strictly isomorphic to $\alpha_q\oplus(c,J_2)$, where $J_2$ is the
indecomposable two-dimensional local system on $c$.
\item There is a term $\kappa_q\in\operatorname{End}(N_q)$, which replaces one arrow of the differential near
the corner $(0,0)$, such that $b+\kappa_q$ is a Maurer--Cartan element and $(N_q,b+\kappa_q)$ is strictly
isomorphic to $\mathrm{BN}_q$.
\end{enumerate}
For every odd $q$ with $5\le q\le21$, $\mathrm{BN}_q$ is an arc obtained from $\alpha_q$ by sliding one of its
ends along~$c$.
\end{theorem}

\begin{theorem}\label{thm:pairing}
Let $5\le q\le21$ be odd, and let $\Gamma$ be a linear curve parallel neither to $\alpha_q$ nor to $c$.
Then the minimal intersection numbers satisfy $i(\Gamma,\mathrm{BN}_q)=i(\Gamma,\alpha_q)+2\,i(\Gamma,c)$,
and if $\Gamma$ is bigraded, with the corner $(\pi,\pi)$ as the special puncture of~\cite{KWZ}, then,
granting the arc case of the intersection formula of Kotelskiy, Watson and Zibrowius~\cite[Thm.~5.25]{KWZ},
$\dim\HF(\Gamma,\mathrm{BN}_q)=\dim\HF\bigl(\Gamma,\alpha_q\oplus(c,J_2)\bigr)$.
\end{theorem}

\begin{corollary}\label{cor:closures}
Granting the arc case of the intersection formula of Kotelskiy, Watson and Zibrowius~\cite[Thm.~5.25]{KWZ}, for
$q\in\{5,7\}$ and every $r\in\Q\cup\{\infty\}$,
\[
\dim\HF\bigl(E_r,(N_q,b)\bigr)=\dim\HF\bigl(E_r,(N_q,b+\kappa_q)\bigr)=\rk\widetilde{\Kh}\bigl(K_r(q);\F_2\bigr).
\]
\end{corollary}

Thus, granting that case, no rational closure distinguishes $(N_q,b)$ from $\mathrm{BN}_q$; curves parallel to $c$ do
(Remark~\ref{rem:sharpness}). Section~\ref{sec:khovanov} states these results in precise form
(Theorems~\ref{thm:structure-full} and~\ref{thm:pairing-full}, Corollary~\ref{cor:closures-full}). The
element $b$ is the switch of a smoothing of the model at one self-intersection, a Lagrangian surgery, and
already reaches the Khovanov rank, whereas $\kappa_q$ is supported on chords around the corner $(0,0)$ and
not at a self-intersection of the curve. Kotelskiy, Watson and Zibrowius suggest that the spectral sequences from
Khovanov homology to Heegaard Floer and instanton homology might be realized by smoothing self-intersections
of $\widetilde{\mathrm{BN}}(T)$ and $\widetilde{\Kh}(T)$~\cite[\S1.8]{KWZ}, and
Theorem~\ref{thm:structure} proceeds in the opposite direction, from a smoothing of the curve on the
instanton side to the Bar-Natan invariant.

\begin{question}\label{q:arc}
Let $q\ge5$ be odd and prime to $3$. If the tangle assignment of Cazassus, Herald, Kirk and Kotelskiy is
defined for $T_q$ with Smith's perturbation for small $t<0$, is the object of $\operatorname{Tw}(\mathcal B)$
that it determines homotopy equivalent to $\mathrm{BN}_q$?
\end{question}

The evidence for $q=5,7$ is Theorem~\ref{thm:structure} and Corollary~\ref{cor:closures}. Two
observations point the other way. First, $\mathrm{BN}_7$ is not homotopy equivalent to any object allowed
by Hypothesis~\ref{hyp:q7-support} (Proposition~\ref{prop:hyp-conflict}). An affirmative answer for
$q=7$ would therefore show that the hypothesis fails, and, granting the arc case of the intersection
formula~\cite[Thm.~5.25]{KWZ}, no rational closure could detect the failure. Second, for $q=11,13$ the encoding of a numerical trace of Smith's curve has pairing dimension
$q+4$ with $E_{-1/2}$, and for $q=17,19$ pairing dimension $q+8$, against $\rk\widetilde{\Kh}(K_q)=q+2$; in
finite searches no deformation by a few switches reaches $\mathrm{BN}_q$ (Remark~\ref{rem:q11}). These
searches do not exclude a larger bounding cochain, and we pose the question without conjecturing an answer.

\subsection{Outline of the arguments}
\emph{The Maurer--Cartan elements at $q=7$.} For small $t$, Smith's perturbed curve contains a component
homeomorphic to the fiber product of the lines of $Q_{1/3}$ and $Q_{1/7}$ with each circle that the Conway sum
acquires over an edge of the pillowcase replaced by two arcs crossing once, joined crosswise to the rest of
the curve (the \emph{cross-resolution} of \S\ref{ss:smith-chart}), and its other components are circles
that can be moved off any fixed complementary curve (Proposition~\ref{prop:smith-main}, after~\cite{Smith}).
Performing the same cross-resolution on the fiber product of the two lines, drawn piecewise linearly, gives a
curve homeomorphic to that component, the model. Encoding the model by its crossings with the arc system of
$\mathcal B$ gives, after unit cancellations, the twisted complex $N=N_7$ on $31$ generators. Morphism
complexes between finite twisted complexes are infinite dimensional, and Lemma~\ref{lem:wrapped-reduction}
computes their homology exactly: the morphism complex is a mapping cone onto free modules over the
polynomial rings of the three faces, and its homology is finite dimensional exactly when those modules are
acyclic over the fraction fields. Each self-intersection of the model has two smoothings; since
$\mathcal B$ has no higher products, the switch $b_s$ of a smoothing that preserves the generators is a
Maurer--Cartan element exactly when $(\delta_N+b_s)^2=0$, a finite check. The enumeration of all $82$
self-intersection orbits (Propositions~\ref{prop:aux-selection} and~\ref{prop:aux-remaining}) compares the
pairing dimensions with the two earrings against the instanton dimensions $9$ and $31$.

\emph{Theorems~\ref{thm:structure} and~\ref{thm:pairing}.} All objects in the statements are reduced
twisted complexes of loop type, and a strict isomorphism between two of them is decided by comparing
the words of their components; $\widetilde{\mathrm{BN}}(T_q)$ is computed with Zibrowius's
program~\cite{Zibkht}. For Theorem~\ref{thm:pairing} we realize each curve by a geodesic in a metric of
nonpositive curvature (the Alexander--Berg--Bishop theorem, see~\cite[Thm.~7.28]{AKP}), in which minimal
intersection numbers are counted locally. A lift of $\mathrm{BN}_q$ agrees with the union of a lift of
$\alpha_q$ and two periods of a lift of $c$ outside small discs at five lattice points, and a four-case
local count at one lattice point~\eqref{eq:local-identity} gives the identity of intersection numbers. For
bigraded curves we then use the intersection formulas of Kotelskiy, Watson and
Zibrowius~\cite[Thm.~5.25]{KWZ} and of Zibrowius~\cite[Thm.~4.45]{Zibrowius}.

\subsection{Scope of claims}
Proposition~\ref{prop:q7-typeD}, Theorems~\ref{thm:structure} and~\ref{thm:pairing} and
Corollary~\ref{cor:closures} are statements about explicitly presented objects of
$\operatorname{Tw}(\mathcal B)$; they neither assume nor prove Conjecture~D of~\cite{CHKK}, and
they do not identify the object that the conjecture of Cazassus, Herald, Kirk and Kotelskiy assigns to $T_q$.
Proposition~\ref{prop:smith-main} identifies the homeomorphism type of the main components of Smith's curve;
it is not proved that their twisted complex is $N_q$. This is corroborated for $q\le13$ by numerical traces
(Remark~\ref{rem:analytic}, \S\ref{ss:methods}), and Hypothesis~\ref{hyp:q7-support} is formulated for the
model. Corollary~\ref{cor:aux-conditional} is conditional on Conjecture~D of~\cite{CHKK} and on
Hypothesis~\ref{hyp:q7-support}. The second statement of Theorem~\ref{thm:pairing}, Corollary~\ref{cor:closures} and the second statement of Proposition~\ref{prop:hyp-conflict} grant the arc case of the intersection formula of Kotelskiy, Watson and
Zibrowius~\cite[Thm.~5.25]{KWZ}, the case of an arc that ends at a puncture, which they assert is treated
like the compact case but do not write out (proof of Theorem~\ref{thm:pairing-full}); the dimensions of
Corollary~\ref{cor:closures} are also computed directly at several hundred slopes. The identifications of $N_q$ and of the switches with the encodings and smoothings of the piecewise-linear curves
rest on the floating-point geometry of these curves (Appendix~\ref{ss:methods}), and so do the second part of
Proposition~\ref{prop:q7-typeD}, Propositions~\ref{prop:aux-selection}--\ref{prop:aux-span},
Corollary~\ref{cor:aux-conditional} and Proposition~\ref{prop:hyp-conflict}, whereas all statements about
the displayed complexes are exact. Statements labelled \emph{Computation} (Appendix~\ref{sec:compute}) are outputs of
finite procedures and are recorded as data.

\subsection{Organization}
Section~\ref{sec:background} recalls the pillowcase, Smith's coordinates and the algebra $\mathcal B$.
Section~\ref{sec:smith} treats Smith's decomposition with zero bounding cochain and the Maurer--Cartan
elements at $q=7$, and Section~\ref{sec:khovanov} relates them to the Bar-Natan invariant.
Section~\ref{sec:disc} discusses related work and open problems, and Appendix~\ref{app:data} prints the
remaining data, records finite polygon counts, and describes the computations.

\subsection*{Acknowledgements}
The pillowcase model and the geometry used here are due to Hedden, Herald and Kirk~\cite{HHK1,HHK2},
Cazassus, Herald, Kirk and Kotelskiy~\cite{CHKK}, Herald and Kirk~\cite{HK}, and Smith~\cite{Smith}. The
immersed-curve invariants are those of Kotelskiy, Watson and Zibrowius~\cite{KWZ}, computed with Zibrowius's
program~\cite{Zibkht}, and the ranks of reduced Khovanov homology were computed with Khoca~\cite{Khoca}.
